\section{Search}
\label{sec:search}

\subsection{Matching and highlighting}

The document's text is not available as a string; it is a set of positioned glyph runs in a
text layer. The find index concatenates each page's run strings into one haystack and records
the run boundaries, so a match is an offset into a page's concatenated text plus a length.

Highlighting does not compute its own geometry. It rewrites the text layer's spans, splitting
a span at the match boundaries and wrapping the remainder in a \texttt{<mark>}-like element,
which is the technique pdf.js uses internally. The benefit is that a painted match is aligned
with the painted glyphs \emph{by construction}: there is no second coordinate computation that
could disagree with the first.

Because the index and the layer are separate, they can disagree. Before painting, the page's
live span strings are compared with the strings the index was built from; a page that does not
match is skipped rather than mis-highlighted. In practice this triggers when a page's text
layer is re-rendered between indexing and searching --- on a fast scroll through a document
whose index is still being built --- and skipping is the correct answer, since the offsets no
longer mean what they did.

\subsection{Result presentation}

Results are grouped by page, in document order, each row quoting the words around the match.
The quote comes from the index itself rather than from re-extracting the page: the offsets a
hit holds are offsets into exactly that string, so the snippet cannot drift from the
highlight.

Two ordering defects are worth recording because both produced a screen that contradicted
itself. First, the panel was told about a query \emph{before} the search ran, so it redrew from
the previous query's matches; because the zero-hit path returns early, nothing redrew
afterwards, and a search for nonsense left twenty-two results beside a status bar saying there
were none. Second, a match list is a view of mutable state: the ``current'' row must be
re-derived whenever the index changes, not only when the user moves to the next match.

\subsection{Folder search across documents}
\label{sec:folder-search}

Searching a folder means reading files the user has not opened. The application's desktop shell
refuses to read any path the user has not granted, so listing a folder also \emph{grants} the
files it finds: a folder the user types is as explicit a permission as one chosen from a
dialog, and the grant is what lets the scan proceed. Depth and file count are bounded, and a
file that will not parse is skipped rather than fatal --- one damaged document in a folder
should not stop the search.

Each file is opened on its own and released again. A folder search is a scan, not a shelf of
open documents; holding every parsed document would cost far more memory than it saves, and
\S\ref{sec:eval} shows the index is cheap enough to rebuild.

\subsection{Saved searches}

A saved search stores the question, not the answer: the query, the case and whole-word flags,
and the folder when there was one. The hits go stale the moment a document changes, and
re-running is the point. Saving under a name that already exists replaces it, so re-saving
updates rather than accumulating near-duplicates, and an entry with no name or no query is
dropped on read rather than repaired.

\subsection{Form fields, without a form layer}
\label{sec:fields}

The viewer highlights form fields --- tinting the rectangles of \texttt{/Widget} annotations so
a reader can see where a form asks for input --- and does not claim to fill them. The
distinction matters for the paper because it is the honest boundary of the feature: detecting
a widget needs the page's annotation array, while filling one needs the field's own hierarchy
(an \texttt{/AcroForm} dictionary, field types, appearance streams, calculation order and
validation), which this application does not implement. Highlighting without filling is
useful and truthful; a half-implemented form layer would be neither.
